\section*{Sets and Indices}

\begin{itemize}
    \item Products $i \in P = \{X,Y\}$.
\end{itemize}

\section*{Parameters}

\begin{itemize}
    \item Available machine time in the week, in hours: $A^{M}$ (code: \texttt{machine\_avail}) $=40$.
    \item Available craftsman time in the week, in hours: $A^{C}$ (code: \texttt{craftsman\_avail}) $=35$.
    \item Machine-time cost, in GBP per hour: $c^{M}$ (code: \texttt{machine\_cost}) $=10$.
    \item Craftsman-time cost, in GBP per hour: $c^{C}$ (code: \texttt{craftsman\_cost}) $=2$.
    \item Revenue per batch of product $i$: $r_i$, with
    \[
    r_X \text{ (code: \texttt{rev\_X}) } = 20,\qquad
    r_Y \text{ (code: \texttt{rev\_Y}) } = 30.
    \]
    \item Minimum required batches of product $X$: $L_X$ (code: \texttt{min\_X}) $=10$.
    \item Machine processing time per batch of product $i$, in minutes: $m_i$, with
    \[
    m_X = 13,\qquad m_Y = 19.
    \]
    \item Craftsman processing time per batch of product $i$, in minutes: $k_i$, with
    \[
    k_X = 20,\qquad k_Y = 29.
    \]
\end{itemize}

\section*{Decision Variables}

\begin{itemize}
    \item Batches of product $X$ to produce: $x_X$ (code: \texttt{X}), continuous, $x_X \ge 0$.
    \item Batches of product $Y$ to produce: $x_Y$ (code: \texttt{Y}), continuous, $x_Y \ge 0$.
\end{itemize}

\section*{Objective}

Maximize weekly profit, defined as revenue minus machine-time and craftsman-time costs:
\[
\max \; z
=
r_X x_X + r_Y x_Y
- c^{M}\left(\frac{m_X x_X + m_Y x_Y}{60}\right)
- c^{C}\left(\frac{k_X x_X + k_Y x_Y}{60}\right).
\]

With the instance data substituted,
\[
\max \; z
=
20x_X + 30x_Y
- 10\left(\frac{13x_X+19x_Y}{60}\right)
- 2\left(\frac{20x_X+29x_Y}{60}\right).
\]

\section*{Constraints}

Machine-time availability:
\[
\frac{m_X x_X + m_Y x_Y}{60} \le A^{M}.
\]

Craftsman-time availability:
\[
\frac{k_X x_X + k_Y x_Y}{60} \le A^{C}.
\]

Minimum contractual production of product $X$:
\[
x_X \ge L_X.
\]

Nonnegativity:
\[
x_X \ge 0,\qquad x_Y \ge 0.
\]

For this instance, the constraints are
\[
\frac{13x_X+19x_Y}{60} \le 40,
\]
\[
\frac{20x_X+29x_Y}{60} \le 35,
\]
\[
x_X \ge 10,
\]
\[
x_X,x_Y \ge 0.
\]

\section*{Notes on Implementation}

\begin{itemize}
    \item The code constructs a linear program with two continuous decision variables using \texttt{LpProblem(..., LpMaximize)}.
    \item Resource usage expressions are formed in hours by dividing the per-batch minute coefficients from \texttt{table\_1.csv} by $60$:
    \[
    \texttt{machine\_hours\_used}=\frac{13\,\texttt{X}+19\,\texttt{Y}}{60},\qquad
    \texttt{craftsman\_hours\_used}=\frac{20\,\texttt{X}+29\,\texttt{Y}}{60}.
    \]
    \item The objective includes only the costs of \emph{used} machine and craftsman hours. Idle time has no explicit variable and no cost term.
    \item No integrality restrictions are imposed; fractional batches are allowed exactly as implemented.
    \item The model is solved by the Gurobi command-line solver through \texttt{GUROBI\_CMD(msg=0)}.
\end{itemize}
